\begin{figure*}[t]
\centering
\scriptsize
\setlength{\fboxsep}{1em}
\setlength{\fboxrule}{1pt}
\fcolorbox{black}{white}{
\parbox{0.955\textwidth}{
\begin{alltt}


You will be provided instructions for completing a task followed by a task to complete.

Instructions:

1. Understand the context: You are evaluating research papers to determine if they should be included in a systematic review. You have access to the title, abstract, authors, and journal of each paper.

2. Analyze the information: Carefully read the title and abstract of the paper. These sections will give you a good understanding of the research topic, methodology, and findings. Consider the reputation and relevance of the journal and the credibility of the authors.

3. Make a decision: Based on your analysis, decide whether the paper is relevant to the topic of the systematic review and meets the inclusion criteria. Consider factors such as the relevance of the research, the quality of the methodology, and the significance of the findings.

4. Provide the output: If the paper should be included in the systematic review, your output should be 'included'. If the paper should not be included, your output should be 'not included'.

5. Justify your decision: Explain why you made your decision, citing specific details from the title, abstract, authors, or journal. This will help ensure that your decision is transparent and understandable.

Title: Why people choose teaching: A scoping review of empirical studies, 2007–2016

Abstract: Who enters teaching and why are questions of immense social and political importance throughout the world. This paper presents a scoping review of empirical studies, published between 2007 and 2016, that addressed influences on the choice of teaching as a career. Seventy articles were analysed descriptively and substantively. Our overview of the nature, extent, and range of research published in these articles highlights that most studies focus on motivations for teaching, with intrinsic and altruistic motivations most commonly identified. We argue that a broader range of theoretical perspectives could add fresh insights to the question of why people choose teaching. (PsycINFO Database Record (c) 2018 APA, all rights reserved)

Journal: Teaching and Teacher Education

Follow the instructions to generate an explanation that reasons towards the correct answer to the task above. End the explanation with the correct answer. [/INST] Explanation: 

The title of the paper "Why people choose teaching: A scoping review of empirical studies, 2007-2016" suggests that the paper is focused on the motivations and factors that influence individuals to choose teaching as a career. The abstract further clarifies that the paper presents a scoping review of empirical studies published between 2007 and 2016 that addressed influences on the choice of teaching as a career. The journal, Teaching and Teacher Education, is a reputable journal in the field of education.

Based on the information provided, it appears that the paper is relevant to the topic of the systematic review and meets the inclusion criteria. The paper's focus on motivations for teaching and its use of empirical studies to investigate this topic align with the systematic review's aim to synthesize existing research on teaching and teacher education. Additionally, the paper's publication in a reputable journal suggests that it has undergone rigorous peer-review and is of high quality.

Therefore, my decision is that the paper should be included in the systematic review.

Output: Included.

Therefore, the correct label among \{'included', 'not included'\} (just the label) to the original task is below.
Label: Included
\end{alltt}
}}
\caption{A complete running example of RAFT (Systematic Review Inclusion).}
\label{fig:raft_subset=systematic_review_inclusion_example}
\end{figure*}
